\documentclass[11pt]{article}
% Shared preamble of the RMC kernel write-ups (% Shared preamble of the RMC kernel write-ups (% Shared preamble of the RMC kernel write-ups (\input{preamble} after \documentclass).
\usepackage[margin=1in]{geometry}
\usepackage{amsmath,amssymb,bm}
\usepackage{graphicx}
\usepackage{booktabs}
\usepackage{hyperref}

\newcommand{\dd}{\,\mathrm{d}}
\newcommand{\Ein}{E}
\newcommand{\Eout}{E'}
\newcommand{\Ecm}{E_{c}}
\newcommand{\mucm}{\mu_{c}}
\newcommand{\mulab}{\mu_0}

% Common closing section: how the verification figures are produced
\newcommand{\verificationmethod}{%
Each figure is produced by \texttt{verification/verify\_laws.py}. For an incident
energy $\Ein$ along $+z$, $10^6$ collisions are sampled with MC/DC's own reaction
sampler (\texttt{sample\_elastic\_scattering}, \texttt{sample\_inelastic\_scattering}
or \texttt{sample\_fission}), and every emitted neutron is histogrammed in
$48 \times 24$ bins of $(\Eout, \mulab)$, with $\mulab$ the lab cosine to $+z$. The
inverted kernel used by RMC is integrated over the same bins (Gauss--Legendre panels;
two-body lines are split where $\mulab^*(\Eout)$ crosses a cosine edge), multiplied by
the number of collisions, and compared with a $\chi^2$ test over bins with at least 5
expected counts (the rest pooled into one bin). The tables list $\chi^2/\mathrm{dof}$,
its $p$-value, the emitted neutrons per collision (sampled vs.\ the integral of the
inverted kernel), and the largest relative change of a bin integral when the
quadrature panels are doubled. Left panels: outgoing-energy marginal; middle:
lab-cosine marginal; right: per-bin $z = (\text{sampled} - \text{inverted}) /
\sqrt{\text{inverted}}$.}
 after \documentclass).
\usepackage[margin=1in]{geometry}
\usepackage{amsmath,amssymb,bm}
\usepackage{graphicx}
\usepackage{booktabs}
\usepackage{hyperref}

\newcommand{\dd}{\,\mathrm{d}}
\newcommand{\Ein}{E}
\newcommand{\Eout}{E'}
\newcommand{\Ecm}{E_{c}}
\newcommand{\mucm}{\mu_{c}}
\newcommand{\mulab}{\mu_0}

% Common closing section: how the verification figures are produced
\newcommand{\verificationmethod}{%
Each figure is produced by \texttt{verification/verify\_laws.py}. For an incident
energy $\Ein$ along $+z$, $10^6$ collisions are sampled with MC/DC's own reaction
sampler (\texttt{sample\_elastic\_scattering}, \texttt{sample\_inelastic\_scattering}
or \texttt{sample\_fission}), and every emitted neutron is histogrammed in
$48 \times 24$ bins of $(\Eout, \mulab)$, with $\mulab$ the lab cosine to $+z$. The
inverted kernel used by RMC is integrated over the same bins (Gauss--Legendre panels;
two-body lines are split where $\mulab^*(\Eout)$ crosses a cosine edge), multiplied by
the number of collisions, and compared with a $\chi^2$ test over bins with at least 5
expected counts (the rest pooled into one bin). The tables list $\chi^2/\mathrm{dof}$,
its $p$-value, the emitted neutrons per collision (sampled vs.\ the integral of the
inverted kernel), and the largest relative change of a bin integral when the
quadrature panels are doubled. Left panels: outgoing-energy marginal; middle:
lab-cosine marginal; right: per-bin $z = (\text{sampled} - \text{inverted}) /
\sqrt{\text{inverted}}$.}
 after \documentclass).
\usepackage[margin=1in]{geometry}
\usepackage{amsmath,amssymb,bm}
\usepackage{graphicx}
\usepackage{booktabs}
\usepackage{hyperref}

\newcommand{\dd}{\,\mathrm{d}}
\newcommand{\Ein}{E}
\newcommand{\Eout}{E'}
\newcommand{\Ecm}{E_{c}}
\newcommand{\mucm}{\mu_{c}}
\newcommand{\mulab}{\mu_0}

% Common closing section: how the verification figures are produced
\newcommand{\verificationmethod}{%
Each figure is produced by \texttt{verification/verify\_laws.py}. For an incident
energy $\Ein$ along $+z$, $10^6$ collisions are sampled with MC/DC's own reaction
sampler (\texttt{sample\_elastic\_scattering}, \texttt{sample\_inelastic\_scattering}
or \texttt{sample\_fission}), and every emitted neutron is histogrammed in
$48 \times 24$ bins of $(\Eout, \mulab)$, with $\mulab$ the lab cosine to $+z$. The
inverted kernel used by RMC is integrated over the same bins (Gauss--Legendre panels;
two-body lines are split where $\mulab^*(\Eout)$ crosses a cosine edge), multiplied by
the number of collisions, and compared with a $\chi^2$ test over bins with at least 5
expected counts (the rest pooled into one bin). The tables list $\chi^2/\mathrm{dof}$,
its $p$-value, the emitted neutrons per collision (sampled vs.\ the integral of the
inverted kernel), and the largest relative change of a bin integral when the
quadrature panels are doubled. Left panels: outgoing-energy marginal; middle:
lab-cosine marginal; right: per-bin $z = (\text{sampled} - \text{inverted}) /
\sqrt{\text{inverted}}$.}

\usepackage{amsthm}

\newtheorem{proposition}{Proposition}
\newcommand{\tpsi}{\tilde\psi}
\newcommand{\teps}{\tilde\epsilon}
\newcommand{\Sbar}{\bar S}
\newcommand{\Sigt}{\Sigma_t}
\newcommand{\E}{\mathbb{E}}

\title{Residual Monte Carlo with a linear discontinuous angular trial space}
\date{}

\begin{document}
\maketitle

\noindent Slab geometry, $\mu=\boldsymbol\Omega\cdot\hat{\mathbf z}$; the angular flux is
integrated over azimuth, $\psi(z,E,\mu)=\int_0^{2\pi}\psi(z,E,\boldsymbol\Omega)\dd\varphi$.
RMC iterates $r=q-L\tpsi$, $L\epsilon=r$ (Monte Carlo), $\tpsi\leftarrow\tpsi+\Pi\epsilon$
[1--5]; see \texttt{continuous\_linear\_space.tex} for the spatial and
\texttt{linear\_energy\_basis.tex} for the energy part of the trial space.

\section{Why linear in angle}

With $\tpsi$ constant on each polar bin $\mu\in[\mu_j,\mu_{j+1}]$, the in-bin angular
shape of the true flux generally leaves an unresolved residual, like the in-bin
energy shape: the streaming term $\mu\,\partial_z\tpsi$ is linear in $\mu$ inside the
bin while $\tpsi$ is not. Streaming can produce angular anisotropy near boundaries,
interfaces and localized sources; isotropic scattering itself still produces an
isotropic scattering source. A linear representation
in $\mu$ reduces the projection error of a smooth angular flux from $O(\Delta\mu)$ to
$O(\Delta\mu^2)$. Linear discontinuous (LD) angular trial spaces were used in the
exponentially convergent Monte Carlo work in slab geometry: in space and direction with
h-adaptivity [2], and with anisotropic discrete-angle scattering, including the
$\mu$-slope moments of the scattering source [3].

\section{Trial space}

On bin $j$ with width $\Delta\mu_j$, let $y_j(\mu)=(2\mu-\mu_j-\mu_{j+1})/\Delta\mu_j$
and
\begin{equation}
  \tpsi(\mu) = \sum_{b=0}^{1}\tpsi_{jb}\,P_b\big(y_j(\mu)\big),\qquad P_0=1,\ P_1=y .
\end{equation}
The grid keeps $\mu=0$ as an edge: the exact angular flux is discontinuous at $\mu=0$ at
a vacuum boundary (incoming directions carry no flux) and its $\mu$-derivative is
singular there in general, so continuity in $\mu$ across $0$ would be wrong. With a
single bin on each half-range the space is the double-$P_1$ ($DP_1$) approximation of
Yvon [6], the half-range expansion suited to slabs with vacuum boundaries [7]; finer
bins give a discontinuous piecewise-linear angular finite element space in slab
angle (the space--direction construction in [2]). The spatial DG references [8]
are background for discontinuous elements, not sources for angular discretization.

\paragraph{Boundary conditions.} Vacuum: $\tpsi_{jb}=0$ on incoming bins of the boundary
trace (the boundary node for continuous spatial elements). An unconstrained spatial
trace requires a boundary residual. Reflective: $\tpsi(-\mu)=\tpsi(\mu)$; for a bin $j$ and its mirror $j^*$,
$y_{j^*}(-\mu)=-y_j(\mu)$, so $\tpsi_{j^*0}=\tpsi_{j0}$ and $\tpsi_{j^*1}=-\tpsi_{j1}$.

\section{Projection and its estimator}

The Legendre basis is orthogonal on each bin, so the $L^2$ projection is local:
\begin{equation}
  (\Pi\psi)_{jb} = \frac{2b+1}{\Delta\mu_j}\int_j \psi(\mu)\,P_b(y_j(\mu))\dd\mu .
\end{equation}
Since $\mu$ is constant along a track, the track-length estimator of the $b=1$ moment
scores $w\,\ell\,y_j(\mu)$ per segment, the angular analogue of the energy slope.
The identity $\E[y^2]=1/3$ under uniform point sampling does not establish a fixed
relative-noise penalty for track-length tallies. Their covariance depends on the
sampled tracks and weights, and relative error is undefined for a zero slope.
History requirements must be assessed for the actual observables.

\begin{proposition}
If $L^{-1}$ exists and the exact full residual, including external-source and
boundary remainders, is transported with an unbiased finite-expectation tally,
$\E[\tpsi^{(n+1)}\mid\tpsi^{(n)}]=\Pi\psi$ for every $\tpsi^{(n)}$ in
the trial space; the angular part of $\Pi$ is the $L^2$ projection onto piecewise-linear
functions of $\mu$ on the bins. This expectation identity alone does not imply
stochastic convergence.
\end{proposition}

\section{Angular transfer}

For an incident direction $\mu'$ and a lab scattering cosine $\mu_0$, with the azimuth
uniform, the outgoing polar cosine is
\begin{equation}
  \mu(\theta) = c + d\cos\theta,\qquad c=\mu'\mu_0,\quad d=\sqrt{1-\mu'^2}\sqrt{1-\mu_0^2},
  \quad\theta\in[0,\pi]\ \text{uniform},
\end{equation}
the substitution that removes the inverse-square-root singularity of the azimuthal
kernel $K(\mu',\mu,\mu_0) = 1/\big(\pi\sqrt{1-\mu'^2-\mu^2-\mu_0^2+2\mu'\mu\mu_0}\big)$
[3]. Let $\theta_{\rm lo}=\arccos\big(\mathrm{clip}((\mu_{j+1}-c)/d)\big)$ and
$\theta_{\rm hi}=\arccos\big(\mathrm{clip}((\mu_j-c)/d)\big)$ be the azimuths where
$\mu(\theta)$ enters and leaves bin $j$. The outgoing moments are then closed form:
\begin{align}
  \Pi^{0}_j(\mu',\mu_0) &= \frac1\pi\int_{\theta_{\rm lo}}^{\theta_{\rm hi}}\dd\theta
  = \frac{\theta_{\rm hi}-\theta_{\rm lo}}{\pi},\\
  \Pi^{1}_j(\mu',\mu_0) &= \frac1\pi\int_{\theta_{\rm lo}}^{\theta_{\rm hi}}
  y_j\big(\mu(\theta)\big)\dd\theta
  = \frac{2}{\pi\Delta\mu_j}\Big[(c-\bar\mu_j)(\theta_{\rm hi}-\theta_{\rm lo})
  + d\,(\sin\theta_{\rm hi}-\sin\theta_{\rm lo})\Big],
\end{align}
with $\bar\mu_j$ the bin center. The piecewise-constant scheme uses only $\Pi^0_j$ (its
arcsine form). The angular transfer between incident bin $j'$ (weighted by
$P_{b'}$) and outgoing bin $j$ (weighted by $P_b$) is
\begin{equation}
  A^{b'b}_{j'j}(\mu_0) = \int_{j'}\dd\mu'\,P_{b'}\big(y_{j'}(\mu')\big)\,
  \Pi^{b}_j(\mu',\mu_0),
\end{equation}
an integral with kinks where $\mu(\theta)$ is tangent to a bin edge
($\mu' = \mu_e\mu_0\pm\sqrt{1-\mu_e^2}\sqrt{1-\mu_0^2}$ for each edge $\mu_e$ of bin
$j$), the same breakpoints as for $b=b'=0$. For $\sum_j$ and $b=0$ the outgoing
probabilities sum to one, so $\sum_jA^{b'0}_{j'j}(\mu_0)=\int_{j'}P_{b'}\dd\mu'$ (which
is $\Delta\mu_{j'}$ for $b'=0$ and $0$ for $b'=1$): a check of any implementation.

The transfer moments of \texttt{rmc\_math.tex} acquire the angular indices,
\begin{equation}
  M_{m,\,a'g'j'b',\,agjb} = \sum_{r\in m}N_r\int_{g'}\dd E'\,\sigma_r(E')P_{a'}
  \int_g\dd E\,P_a\int\dd\mu_0\,f_r(E'\to E,\mu_0)\,A^{b'b}_{j'j}(\mu_0),
\end{equation}
and the projected in-scatter source is
$\Sbar_{gjab}=\frac{(2a+1)(2b+1)}{\Delta E_g\Delta\mu_j}
\sum_{a'g'j'b'}M_{m,\,a'g'j'b',\,agjb}\tpsi_{g'j'a'b'}$, where $N_r$ is the
nuclide's atom density. With linear
energy and angle the moment arrays are $16\times$ those of the piecewise-constant
scheme; the kernel evaluations are the same.

\paragraph{Isotropic and linearly anisotropic scattering.} For isotropic lab emission,
the scattering-cosine average, not the transfer at a fixed $\mu_0$, is
\[
  \overline A^{b'b}_{j'j}=\frac12\int_{-1}^1 A^{b'b}_{j'j}(\mu_0)\dd\mu_0
  =\frac12\left(\int_{j'}P_{b'}\dd\mu'\right)
             \left(\int_jP_b\dd\mu\right).
\]
This vanishes unless $b=b'=0$: the bin slopes do not contribute to an isotropic
source. At fixed $\mu_0=1$, by contrast, outgoing and incident directions coincide,
and $A^{11}_{jj}(1)=\Delta\mu_j/3$. For a $P_1$ kernel
$1+3\bar\mu_0\mu_0$ (per $4\pi$), the
addition theorem gives the outgoing first moment
$\int\dd\mu\,\mu\,(\ldots)\propto\bar\mu_0\,\mu'$, so only the incident flux's first
angular moment transfers anisotropy. These limits check $A^{b'b}$ against closed forms.

\section{Residual}

With $\tpsi$ linear in $\mu$ on each bin,
\begin{equation}
  r_c = Q + \Sbar - \Sigt\tpsi - \mu\,\partial_z\tpsi ,
\end{equation}
with the full residual also including $r_s$, any fission remainder and $q-Q$.
The collision terms are linear in $\mu$ on the bin and the streaming term is
\emph{quadratic} in $\mu$ (the factor $\mu$ times a linear function). With a continuous
linear spatial space, $\partial_z\tpsi$ is constant on each cell, so for fixed $E$ the
residual on a bin $(k,j)$ is a polynomial of degree one in $s$ and two in $\mu$.
Writing it as $A(\mu)+sB(\mu)$, the integral over $s\in[0,1]$ is
$|A+B/2|$ if the endpoint values have the same sign, and
$[A^2+(A+B)^2]/(2|B|)$ if they have opposite signs. For $B=0$ it is $|A|$.
The resulting angular integrand can be rational; splitting at its sign changes
does not make Gaussian quadrature exact. Use analytic integration where available
or numerical quadrature with an error estimate. Sampling is as before (bin by mass,
uniform inside, weight
$r_c/q$); boundary faces carry no residual with the strong boundary conditions above.
With a piecewise-constant or discontinuous spatial space, the face residual
$-\mu\,[\tpsi]\,\delta(z-z_f)$ has density $\propto|\mu\,(D_0+D_1y)|$ on the bin; sampling
$\mu\propto|\mu|$ and weighting by $(D_0+D_1y)\,\mathrm{sign}(-\mu)$ is unbiased.

The scattering correction $r_s=\mathcal S\tpsi-\Sbar$ is unchanged in form: the
in-scatter density integrates the incident flux, now linear in $\mu'$ on each bin,
against the emission kernel, which needs the kernel's bin moments
$\int_{j'}P_{b'}(y_{j'}(\mu'))\,\tau(\mu'\to\mu)\dd\mu'$.
The angular bin integrals of $K$ and $P_1K$ are closed form in the substitution
above; a continuous emission law still requires integration over its scattering
cosine and incident energy.

\section{Beyond the slab}

In 2D and 3D the polar cosine about one axis is replaced by the direction on the sphere.
The two standard low-order choices are spherical harmonics ($P_1$: $\psi\approx(\phi+3
\mathbf J\cdot\boldsymbol\Omega)/4\pi$), for which the angular transfer reduces to the
Legendre moments of the scattering kernel by the addition theorem [7], and discontinuous
piecewise-linear functions on a partition of the sphere. Angular finite elements
on spherical geodesic grids are studied in [9]; a discontinuous extension is a
separate design choice.
In RMC unresolved angular residuals contribute to finite-history noise, whose
size also depends on the proposal and tally response. Vermaak and Morel used an
isotropic approximate solution and
still found RMC more efficient than standard Monte Carlo in their continuous-space
benchmarks [4]. A global isotropic or $P_1$ function cannot satisfy vacuum on an
entire incoming hemisphere without also vanishing on outgoing directions. Retain
the boundary residual, as [4] does, or choose a compatible half-range space.

\begin{thebibliography}{99}
\bibitem{halton} J.~H.~Halton, ``Sequential Monte Carlo,'' \emph{Proc.\ Cambridge
  Philos.\ Soc.} \textbf{58}, 57 (1962).
\bibitem{peterson} J.~R.~Peterson, J.~E.~Morel and J.~C.~Ragusa,
  ``Exponentially-convergent Monte Carlo for the 1-D transport equation,'' M\&C 2013;
  J.~R.~Peterson, M.S.\ thesis, Texas A\&M University (May 2014),
  \href{https://oaktrust.library.tamu.edu/bitstreams/69ab2dae-3fef-4daf-96d0-2bd1896d81e3/download}
  {university copy}.
  Published formulation: \emph{Residual Monte Carlo for the One-Dimensional
  Particle Transport Equation}, SIAM J.\ Sci.\ Comput.\ \textbf{38}(6),
  B941--B961 (2016), \href{https://doi.org/10.1137/16M1057619}{doi:10.1137/16M1057619}.
\bibitem{franke} B.~C.~Franke, D.~E.~Bruss and J.~E.~Morel, ``Residual Monte Carlo with
  Discrete Scattering Angles in the 1-D Transport Equation,'' ANS Winter Meeting
  \emph{Trans.\ Am.\ Nucl.\ Soc.} \textbf{109}(1), 679--682 (2013),
  \href{https://www.ans.org/pubs/transactions/volume-109/}{ANS publication record}.
\bibitem{vermaak} J.~I.~C.~Vermaak and J.~E.~Morel, ``Transport Error Estimation Using
  Residual Monte Carlo,'' \emph{J.\ Comput.\ Phys.} (2022);
  https://www.sciencedirect.com/science/article/pii/S0021999122003680, OSTI 1977284.
\bibitem{larsen} M.~Larsen et al., ``A Residual Monte Carlo Algorithm for Continuous
  Energy Neutron Transport with Elastic Scattering,'' \emph{Nucl.\ Sci.\ Eng.}
  \textbf{200}(3), 525--538 (2026), doi:10.1080/00295639.2025.2495608;
  ``A Residual Monte Carlo Algorithm for 1D Continuous-Energy
  Neutron Transport,'' \emph{Nucl.\ Sci.\ Eng.} (2026), doi:10.1080/00295639.2026.2659992.
\bibitem{yvon} J.~Yvon, ``La diffusion macroscopique des neutrons: une m\'ethode
  d'approximation,'' \emph{J.\ Nucl.\ Energy} \textbf{4}, 305 (1957).
\bibitem{lewis} E.~E.~Lewis and W.~F.~Miller, \emph{Computational Methods of Neutron
  Transport}, Wiley (1984) (half-range and double-$P_N$ expansions, spherical
  harmonics, the addition theorem).
\bibitem{reed} W.~H.~Reed and T.~R.~Hill, ``Triangular mesh methods for the neutron
  transport equation,'' LA-UR-73-479, Los Alamos (1973); P.~Lesaint and P.~A.~Raviart,
  ``On a finite element method for solving the neutron transport equation,'' in
  \emph{Mathematical Aspects of Finite Elements in Partial Differential Equations},
  Academic Press (1974).
\bibitem{angular-fe} M.~K.~Bhattacharyya and D.~Radice,
  \emph{A Finite Element Method for Angular Discretization of the Radiation
  Transport Equation on Spherical Geodesic Grids}, J.\ Comput.\ Phys.\ (2023),
  \href{https://doi.org/10.1016/j.jcp.2023.112365}{doi:10.1016/j.jcp.2023.112365}.
\end{thebibliography}

\end{document}
